\documentclass[dvipdfmx]{beamer}
%テーマの設定
\usetheme[numbering=fullbar]{focus}

%パッケージの設定
\usepackage{graphicx}
\usepackage{pxjahyper}
\usepackage{tikz}
\usetikzlibrary{arrows.meta,positioning}
\usepackage{listings, jlisting}
\usepackage{listings-arduino, listings-javascript}
\usepackage{xcolor}

%色設定
\definecolor{identifier}{RGB}{225, 123, 52}       %#e17b34
\definecolor{comment}{RGB}{105, 130, 27}          %#69821b
\definecolor{keyword}{RGB}{170, 76, 143}          %#aa4c8f
\definecolor{string}{RGB}{39, 146, 195}           %#2792c3

%ソースコード設定
\lstset{
  basicstyle={\ttfamily},                         % 基礎の文字のフォント設定
  identifierstyle={\color{identifier}},           % 変数名などのフォント設定
  commentstyle   ={\color{comment}},              % コメントのフォント設定
  keywordstyle   =[1]{\color{keyword}},           % 予約語のフォント設定
  stringstyle    ={\ttfamily\color{string}},      % 文字列のフォント設定
  tabsize=2,
  columns=fixed,
}


%フォント設定
\renewcommand{\kanjifamilydefault}{\gtdefault}

%メタ設定
\title{多言語プログラミング}
\subtitle{プログラミングにおける骨格}
\author[電研]{電気技術研究会}
\institute[NITNC]{NIT, Nara Collage}
\date[\today]{\today}
\titlegraphic{
  \begin{tikzpicture}[overlay, remember picture]
    \node[at=(current page.north east), anchor=north east]{
      \includegraphics[height=0.15\textheight]{pic/Logo(2).png}
      \hspace{0.05\textwidth}
    };
  \end{tikzpicture}
}

\begin{document}

\begin{frame}[plain]
  \titlepage
\end{frame}


\section{概要}

  \begin{frame}{講義目的と進め方}
    講義目的
    \begin{itemize}
      \setlength{\itemsep}{3mm}
      \item \structure{多言語}にわたるプログラミングを\structure{独学する}術を身に付ける
      \item 基本的な\structure{用語}や\structure{技術}を学び、\structure{活かせる力}を身に付ける
      \item マイコンやウェブなど、\structure{広い範囲に適応できる}人材育成
    \end{itemize}
    \vfill
    進め方    
    \begin{itemize}
      \setlength{\itemsep}{3mm}
      \item 重要な\structure{章}ごとにまとめて講義
      \item １日（\structure{２～３時間}）\structure{１～２}章分を進めたい
      \item 各章ごとに\structure{演習}がある
    \end{itemize}
  \end{frame}

  \begin{frame}{カリキュラム}
    \begin{tikzpicture}
      [every node/.style={rounded corners,inner sep=5pt, minimum width=100pt,minimum height=35pt},
      pre-subj/.style={fill=blue!20},post-subj/.style={fill=black!20},overlay, remember picture]
      \node[pre-subj] (a) at (0.15\textwidth,0.375\textheight){プログラミング};
      \node[pre-subj,below=10pt of a] (b) {変数とデータ型};
      \node[pre-subj,below=10pt of b] (c) {演算と演算子};
      \node[pre-subj,below=10pt of c] (d) {制御構文};
      \node[pre-subj,below=10pt of d] (e) {配列};
      \node[pre-subj,right=10pt of e] (f) {オブジェクト};
      \node[pre-subj,above=10pt of f] (g) {関数};
      \node[pre-subj,above=10pt of g] (h) {ポインタ};
      \node[pre-subj,above=10pt of h,align=center] (i) {オブジェクト\\指向構文};
      \node[pre-subj,above=10pt of i] (j) {コンピュータ基礎};
      \node[pre-subj,right=10pt of j,align=center] (k) {マイコン\\プログラミング};
      \node[pre-subj,below=10pt of k] (l) {IoT};
      \node[pre-subj,below=10pt of l,align=center] (m) {プログラミング\\解析};
      \node[pre-subj,below=10pt of m] (n) {コードの可読性};
      \node[pre-subj,below=10pt of n] (o) {言語の学習法};
      \foreach \x / \y in {a/b,b/c,c/d,d/e,e/f,f/g,g/h,h/i,i/j,j/k,k/l,l/m,m/n,n/o}{
        \draw[->,>={Stealth[round]},line width=.5mm] (\x) -- (\y);
      }
    \end{tikzpicture}
  \end{frame}

\section{プログラミングとは}
  \begin{frame}{進度}
    \begin{tikzpicture}
      [every node/.style={rounded corners,inner sep=5pt, minimum width=100pt,minimum height=35pt},
      pre-subj/.style={fill=blue!20},now-subj/.style={white,fill=black!50!green},post-subj/.style={fill=black!20},overlay, remember picture]
      \node[now-subj] (a) at (0.15\textwidth,0.375\textheight){プログラミング};
      \node[pre-subj,below=10pt of a] (b) {変数とデータ型};
      \node[pre-subj,below=10pt of b] (c) {演算と演算子};
      \node[pre-subj,below=10pt of c] (d) {制御構文};
      \node[pre-subj,below=10pt of d] (e) {配列};
      \node[pre-subj,right=10pt of e] (f) {オブジェクト};
      \node[pre-subj,above=10pt of f] (g) {関数};
      \node[pre-subj,above=10pt of g] (h) {ポインタ};
      \node[pre-subj,above=10pt of h,align=center] (i) {オブジェクト\\指向構文};
      \node[pre-subj,above=10pt of i] (j) {コンピュータ基礎};
      \node[pre-subj,right=10pt of j,align=center] (k) {マイコン\\プログラミング};
      \node[pre-subj,below=10pt of k] (l) {IoT};
      \node[pre-subj,below=10pt of l,align=center] (m) {プログラミング\\解析};
      \node[pre-subj,below=10pt of m] (n) {コードの可読性};
      \node[pre-subj,below=10pt of n] (o) {言語の学習法};
      \foreach \x / \y in {a/b,b/c,c/d,d/e,e/f,f/g,g/h,h/i,i/j,j/k,k/l,l/m,m/n,n/o}{
        \draw[->,>={Stealth[round]},line width=.5mm] (\x) -- (\y);
      }
    \end{tikzpicture}
  \end{frame}

  \begin{frame}{プログラミングについて}
    プログラミングどんなもの？
    \begin{itemize}
      \setlength{\itemsep}{3mm}
      \item プログラミングには様々な種類が存在\\ウェブ、マイコン、ゲーム...etc
      \item 大まかには２種類\\\structure{オブジェクト指向}と\structure{構造化プログラミング}（後述）
      \item \structure{骨格}は基本的に共通\\関数や変数など、基本的な\structure{骨格}がある
    \end{itemize}
    \vfill
    プログラミング言語
    \begin{itemize}
      \setlength{\itemsep}{3mm}
      \item \structure{言語}は多数ある\par C, C++, C\#, JavaScript, VisualBasic, Python, \LaTeX
      \item 本教材では主に\structure{JavsScript}とArduinoによる\structure{C++}を使用
    \end{itemize}
  \end{frame}

  \begin{frame}[fragile]{基本的な骨格とは}
    \begin{columns}
      \begin{column}{0.6\textwidth}
        \begin{itemize}
          \item 冒頭で\structure{ライブラリ}の宣言
          \item \structure{メインループ}で処理\\UI等の場合では\structure{イベント}を処理
          \item \structure{サブルーチン（関数）}にて\\繰り返し処理
        \end{itemize}
        \begin{exampleblock}{C言語の例}
          \begin{lstlisting}[language=C]
#include <stdio.h>
int main(void){
  printf('Hello World');
  return 0;
}
    
          \end{lstlisting}
        \end{exampleblock}
      \end{column}
      \begin{column}{0.4\textwidth}
        \begin{exampleblock}{Arduinoの例}
          \begin{lstlisting}[language=Arduino,basicstyle=\footnotesize]
void setup(){
  pinMode(2, OUTPUT);
}
void loop(){
  digitalWrite(2, HIGH);
  delay(200);
  digitalWrite(2,LOW);
  delay(200);
}

          \end{lstlisting}
        \end{exampleblock}
      \end{column}
    \end{columns}
  \end{frame}
  
\end{document}